\documentclass[10pt]{article}

% ---------- Document Identity (update these only) ----------
\newcommand{\DocStage}{}
\newcommand{\DocTitle}{Your Road to Tier One SOF}
\newcommand{\ProgramName}{182-Week Civilian-to-Selection Readiness Pipeline}
\newcommand{\ProgramSubtitle}{}

% ---------- Page ----------
\usepackage[legalpaper,left=0.5in,right=0.5in,top=0.6in,bottom=0.45in]{geometry}

% ---------- Typography ----------
\usepackage[T1]{fontenc}
\usepackage[utf8]{inputenc}
\usepackage{libertinus}
\usepackage{microtype}
\usepackage{parskip}
\usepackage{ragged2e}
\usepackage{textcomp}
\AtBeginDocument{\color{black}}
\AtBeginDocument{\pagecolor{white}}
\AtBeginDocument{\justifying}

% ---------- Landscape pages ----------
\usepackage{pdflscape}

% ---------- Color ----------
\usepackage[table]{xcolor}
\definecolor{StageBlue}{HTML}{1F4E79}
\definecolor{StageBlueDark}{HTML}{163A5A}
\definecolor{LightGray}{HTML}{F2F4F7}
\definecolor{MidGray}{HTML}{D9DEE5}
\definecolor{RuleGray}{HTML}{C7CED8}
\definecolor{HeaderInk}{HTML}{000000}
\definecolor{Ink}{HTML}{000000}

% ---------- Utility ----------
\usepackage{enumitem}
\sloppy
\setlength{\emergencystretch}{2em}

% ---------- Boxes / UI ----------
\usepackage[most]{tcolorbox}

\tcbset{
  charterTitle/.style={
    enhanced,
    colback=StageBlue!4,
    colframe=StageBlue,
    boxrule=1.25pt,
    arc=3pt,
    left=16pt,right=16pt,top=14pt,bottom=14pt,
    borderline north={3.2pt}{0pt}{StageBlue},
    borderline west={4pt}{0pt}{StageBlue},
    borderline south={1.25pt}{0pt}{StageBlue},
    drop shadow={black!25!white},
  },
  charterRule/.style={
    enhanced,
    colback=LightGray,
    colframe=RuleGray,
    boxrule=0.85pt,
    arc=3pt,
    left=12pt,right=12pt,top=10pt,bottom=10pt,
    borderline west={3pt}{0pt}{StageBlue},
  },
  charterBadge/.style={
    enhanced,
    colback=StageBlue,
    coltext=white,
    colframe=StageBlueDark,
    boxrule=0.6pt,
    arc=2pt,
    left=6pt,right=6pt,top=3pt,bottom=3pt,
  }
}
\newtcbox{\Badge}{nobeforeafter,charterBadge}

% ---------- Lists ----------
\setlist[itemize]{leftmargin=1.5em, itemsep=0.25em, topsep=0.25em}
\setlist[enumerate]{leftmargin=1.8em, itemsep=0.25em, topsep=0.25em}

% ---------- TOC ----------
\usepackage{tocloft}
\setcounter{tocdepth}{3}
\setlength{\cftbeforesecskip}{3pt}
\setlength{\cftbeforesubsecskip}{2pt}
\setlength{\cftbeforesubsubsecskip}{0pt}
\renewcommand{\cftsecfont}{\bfseries}
\renewcommand{\cftsecpagefont}{\bfseries}
\renewcommand{\cftsubsecfont}{\normalfont}
\renewcommand{\cftsubsecpagefont}{\normalfont}
\renewcommand{\cftsubsubsecfont}{\normalfont}
\renewcommand{\cftsubsubsecpagefont}{\normalfont}
\setlength{\cftsecindent}{0pt}
\setlength{\cftsubsecindent}{1.25em}
\setlength{\cftsubsubsecindent}{2.8em}
\setlength{\cftsecnumwidth}{2.1em}
\setlength{\cftsubsecnumwidth}{2.8em}
\setlength{\cftsubsubsecnumwidth}{3.6em}

% Remove the built-in TOC heading so only the custom heading is shown
\makeatletter
\renewcommand{\tableofcontents}{\@starttoc{toc}}
\makeatother

% ---------- Header/Footer: page number only ----------
\usepackage{fancyhdr}
\pagestyle{fancy}
\fancyhf{}
\renewcommand{\headrulewidth}{0pt}
\renewcommand{\footrulewidth}{0pt}
\fancyfoot[C]{\thepage}

% ---------- Section formatting ----------
\usepackage{titlesec}

\titleformat{\section}
  {\Large\bfseries\color{Ink}}
  {\thesection}{0.6em}{}
  [\vspace{0.25em}\textcolor{StageBlue}{\rule{\linewidth}{0.8pt}}]

\titleformat{\subsection}
  {\large\bfseries\color{Ink}}
  {\thesubsection}{0.6em}{}

\titleformat{\subsubsection}
  {\normalsize\bfseries\color{Ink}}
  {\thesubsubsection}{0.6em}{}

\titlespacing*{\section}{0pt}{0.95em}{0.35em}
\titlespacing*{\subsection}{0pt}{0.65em}{0.25em}
\titlespacing*{\subsubsection}{0pt}{0.55em}{0.2em}

% ---------- Table engine ----------
\usepackage{tabularray}
\UseTblrLibrary{booktabs}

% ---------- tabularray: GLOBAL table treatment ----------
\SetTblrInner{
  cells = {fg=black, bg=white, valign=t},
}

% ---------- Header helpers ----------
\newcommand{\Hdr}[1]{\shortstack[c]{\strut #1\strut}}

% ---------- Lines ----------
\newcommand{\TitleLine}{\textcolor{StageBlue}{\rule{0.98\linewidth}{0.95pt}}}
\newcommand{\SubTitleLine}{\textcolor{RuleGray}{\rule{0.98\linewidth}{0.65pt}}}

% ---------- Continued on next page ----------
\newcommand{\ContinuedOnNextPageInline}{%
  \par
  \vfill
  \noindent\textcolor{RuleGray}{\rule{\linewidth}{1pt}}\vspace{-2pt}
  \noindent\hfill\textit{Continued on next page}\par\vspace{-10pt}
  \noindent\textcolor{RuleGray}{\rule{\linewidth}{1pt}}
  \clearpage
}

% ---------- PDF hyperlinks ----------
\usepackage[hidelinks]{hyperref}
\hypersetup{
  colorlinks=false,
  pdfborder={0 0 0},
  linktoc=all,
  pdftitle={\DocTitle},
  pdfauthor={\ProgramName}
}

% =========================================================
\begin{document}

\section*{Stage 4.A — Phase 00 Overview — Phase Row Template}
\phantomsection
\addcontentsline{toc}{section}{Stage 4.A — Phase 00 Overview — Phase Row Template}

\subsection*{Layer 1 — Executive Phase Row Summary}
\phantomsection
\addcontentsline{toc}{subsection}{Layer 1 — Executive Phase Row Summary}

\begingroup
\scriptsize
\setlength{\tabcolsep}{2pt}
\renewcommand{\arraystretch}{1.12}

\begin{longtblr}{
  width=\linewidth,
  rowhead=1,
  colspec = {
    Q[l,wd=1.75in]
    X[l,1.00]
  },
  hlines,
  vlines,
  cells = {hsep=6pt, vsep=6pt, halign=l, valign=t},
  row{odd} = {bg=white},
  row{even} = {bg=LightGray},
  row{1} = {bg=StageBlue!15, fg=black, font=\bfseries, halign=l, valign=t},
}
\Hdr{Field} & \Hdr{Entry} \\
Phase \# / Name & Phase 00 — Bodily Reactivation and Baseline Creation \\
Duration (weeks) & 16 \\
Mission & Establish safe movement tolerance and execution discipline under start-from-zero constraints, then prove controlled baseline capacity using \textbf{VALID}, lock-field-governed evidence in Weeks 4, 8, 12, and 16 without introducing high-risk exposures. Phase 00 exists to make the trainee executable and auditable under the 6-session cadence, not to create peak performance outputs. \\
Primary Physiological Focus & neuromuscular re-education; aerobic base via low-impact locomotion; tissue tolerance; movement quality; habit formation; modest initial weight-loss trend; foot-interface stabilization; recovery transparency and readiness signal capture as a safety control \\
Key Training Modalities & low-impact Cardio; calisthenics regressions; mobility and movement prep; breathing downregulation; basic trunk endurance; chair and incline variants; optional light band work only if risk-neutral and does not create a gear dependency; conservative indoor substitutes when environment is unsafe \\
Weekly Structure & High-level 6-session microcycle with Cardio every session and mandatory session structure elements per Stage 1.F, executed as one unified daily session with continuity and complete same-day logging. Cardio intensity is anchored by RPE 0–10 plus talk-test cue; no tier model is required in Phase 00 evidence fields. \\
Progression Model & RPE and talk-test anchored; time-on-feet and tolerance progression; strict lock-field stabilization for any testable items; technique and symptom stability as progression governors; one progression lever at a time; regression expected and acceptable when durability or recovery signals indicate elevated risk \\
Deload Plan & Weeks 4, 8, 12, 16 are Deload+Test windows, hard-locked per Stage 3. These windows are the only authorized weeks for Phase 00 test batteries and for producing gate-grade evidence for Phase 00 exit. \\
Mid-Phase Tests & Week 8 Deload+Test battery: \textbf{SMB} (Requires Stage 8 review - no reconciled Stage 2 \textbf{ID} assigned); \textbf{MOB} Bundle (Requires Stage 8 review - no reconciled Stage 2 \textbf{ID} assigned); \textbf{ME-L-STS-30} (Requires Stage 8 review - no reconciled Stage 2 \textbf{ID} assigned); \textbf{ME-UP-PU-INC} (Requires Stage 8 review - no reconciled Stage 2 \textbf{ID} assigned); \textbf{CAL-PLANK} (Requires Stage 8 review - no reconciled Stage 2 \textbf{ID} assigned); \textbf{AER-WALK-015} (Requires Stage 8 review - no reconciled Stage 2 \textbf{ID} assigned) (mode logged). Week 8 is a confirmation/check-in window; it is not the Phase 00 exit gate. \\
End-Phase Tests & Week 16 Deload+Test exit gate battery: \textbf{SMB} use Stage 2 \textbf{ID}; \textbf{MOB} Bundle use Stage 2 \textbf{ID}; \textbf{ME-L-STS-30} use Stage 2 \textbf{ID}; \textbf{ME-UP-PU-INC} use Stage 2 \textbf{ID}; \textbf{CAL-PLANK} use Stage 2 \textbf{ID}; \textbf{AER-WALK-015} use Stage 2 \textbf{ID} (mode logged); \textbf{REC-ADH-WK} (Requires Stage 8 review - no reconciled Stage 2 \textbf{ID} assigned); \textbf{BODYCOMP} trend directionality review (Requires Stage 8 review - no reconciled Stage 2 \textbf{ID} assigned). Week 16 is the only Phase 00 exit-gate week. \\
Exit Standards & Objective pass/fail criteria tied to \textbf{VALID} evidence only, including Phase 00 numeric gates for adherence, walk time, plank, sit-to-stand, incline push-up, and required mobility elements; no invented thresholds beyond those explicitly listed in this Phase Row. Any \textbf{INVALID} test evidence is treated as not yet tested for gate decisions until repeated in a valid window. \\
Equipment Introduced/Required & Minimal Phase 00 kit only, function-linked; no shopping guidance; dependencies validated before lock-field testing is treated as admissible. Minimum infrastructure emphasizes: footwear interface, timer method, logging method, safe route or indoor substitute, stable chair height and incline height for lock-field tests, hydration access, and basic bodyweight/tape capture capability. \\
Risk Controls & Phase 00-specific conservative posture: low-impact bias; no ruck/run/swim/underwater time trials; no maximal strength; strict stop posture; foot integrity governance; strict validity tagging; environment hazard reslot discipline; lock-field enforcement for testable items; Compromised treated as \textbf{INVALID} for gate use; no “push through” posture is admissible evidence. \\
\end{longtblr}

\endgroup

\newpage

\subsection*{Layer 2 — Full Phase Row Template for Phase 00}
\phantomsection
\addcontentsline{toc}{subsection}{Layer 2 — Full Phase Row Template for Phase 00}

\subsubsection*{1. Phase Mission}
\phantomsection

Phase 00 is the mandatory on-ramp that converts a high-risk sedentary start-state into a controlled, auditable execution posture capable of safely tolerating 6 sessions per week with complete session structure and consistent Cardio. The phase exists to establish baseline movement capacity, tissue tolerance, and low-impact aerobic consistency while enforcing strict governance: complete logging, explicit lock fields, conservative substitutions, and zero tolerance for undocumented drift.

Phase 00 completion is proven by \textbf{VALID} exit-gate evidence in Week 16, produced only inside hard-locked Deload+Test windows (Weeks 4, 8, 12, 16). Phase 00 does not pursue performance peaks and does not introduce high-risk tests. It builds readiness to enter Phase 01 without preventable injury risk escalation or evidence contamination.

Mission boundaries (explicit; Phase 00 is not permitted to drift):

\begin{itemize}
\item Phase 00 does not “prove readiness” through maximal outputs.
\item Phase 00 does not introduce impact running escalation, ruck time trials, swim time trials, underwater work, or maximal strength testing.
\item Phase 00 prioritizes stable execution, stable comparability, and durable tolerance without mechanics breakdown.
\end{itemize}

\subsubsection*{2. Phase Purpose}
\phantomsection

\begin{itemize}
\item Build the minimum viable execution discipline required for the full pipeline:
  \begin{itemize}
  \item stable 6 sessions/week cadence,
  \item complete session structure (Safety Self-Check → warm-up → mobility → main work → Cardio → cooldown → recovery → post-session notes),
  \item same-day logging with no memory backfill,
  \item enforce “no logs = no evidence = no gate credit” posture.
  \end{itemize}

\item Establish a low-impact aerobic base using walk-based instruments and conservative Cardio modalities that minimize orthopedic risk while supporting adherence and body composition directionality.

\item Establish movement-quality baselines using lock-field-governed screening and mobility elements to prevent early “too much too soon” failures and to make regression decisions evidence-driven rather than subjective.

\item Establish Phase 00 measurement comparability:
  \begin{itemize}
  \item lock fields are defined and held stable across Weeks 4–12 for any testable item,
  \item environment and tool identity is stable enough for repeat tests to be comparable,
  \item compromised conditions are treated as \textbf{INVALID} for gate use and are reslotted rather than forced.
  \end{itemize}

\item Establish safety transparency:
  \begin{itemize}
  \item pain is logged with location and 0–10 severity,
  \item mechanics change is treated as a termination trigger for the provoking exposure,
  \item foot hotspot and blister signals are treated as early disqualifiers for impact/load escalation.
  \end{itemize}

\item Produce Week 16 exit-gate evidence that proves safe tolerance and basic capacity for Phase 01 entry without introducing ruck time trials, run time trials beyond walk-based instruments, swim time trials, underwater, or maximal strength tests.
\end{itemize}

\subsubsection*{3. Entry Criteria}
\phantomsection

Phase 00 entry is permitted only when the minimum infrastructure for safe execution and auditable evidence exists.

Minimum entry criteria (governance-level; not a programming prescription):

\begin{itemize}
\item Safety posture ready:
  \begin{itemize}
  \item \textbf{STOP} \textbf{NOW} triggers are understood operationally and enforceable in-session (\textbf{DEP-1C-002}).
  \item Medical escalation pathway exists (\textbf{DEP-1C-002}) and medical clearance posture is satisfied prior to starting training (\textbf{DEP-1C-001}).
  \item Water governance acknowledged: underwater and breath-hold are prohibited in Phase 00; no pool requirement exists for Phase 00 (Stage 0.5 and \textbf{DEP-1C-013} posture).
  \end{itemize}

\item Execution infrastructure ready:
  \begin{itemize}
  \item Single primary logging method exists and is usable daily (\textbf{DEP-1C-005}); no backfilled memory outcomes permitted.
  \item Timing method exists (\textbf{DEP-1C-006}) to verify Cardio continuity and to support time-based monitors; timing method identity is loggable.
  \item Binder access and version visibility exists (\textbf{DEP-1C-015}).
  \end{itemize}

\item Foot interface ready:
  \begin{itemize}
  \item Footwear interface established (\textbf{DEP-1C-007}) and treated as a locked interface for Phase 00 test windows.
  \item Footcare supplies and hotspot control protocol available (\textbf{DEP-1C-008}).
  \item Drying/hygiene posture adequate to avoid predictable maceration breakdown (\textbf{DEP-1C-011}).
  \end{itemize}

\item Environment feasibility ready:
  \begin{itemize}
  \item Safe repeatable walking route OR calibrated treadmill/indoor substitute exists (\textbf{DEP-1C-003}) with environment identity posture that can remain stable for comparability.
  \item Weather and environment substitution plan exists (\textbf{DEP-1C-004}) and is applied as default under unsafe outdoor conditions; unsafe conditions trigger indoor substitution or reslot, not forced compliance.
  \end{itemize}

\item Schedule stability ready:
  \begin{itemize}
  \item Protected time blocks exist to support 6 sessions per week (\textbf{DEP-1C-014}).
  \end{itemize}

\item Phase 00 test lock-field readiness by Week 4 (required for gate-grade validity):
  \begin{itemize}
  \item Chair, bench, or step capable of locked height for \textbf{ME-L-STS-30} (Requires Stage 8 review - no reconciled Stage 2 \textbf{ID} assigned).
  \item Incline surface capable of locked height for \textbf{ME-UP-PU-INC} (Requires Stage 8 review - no reconciled Stage 2 \textbf{ID} assigned).
  \end{itemize}
\end{itemize}

If any entry criterion is not met:

\begin{itemize}
\item Default gate outcome for Phase start is \textbf{HOLD} (or \textbf{MEDICAL} \textbf{HOLD} when stop-rule triggers or clinician restrictions apply).
\item Exposure is regressed to the safest feasible execution until dependencies are satisfied.
\item Planned Week 4 testing is tagged \textbf{INVALID} if lock fields cannot be satisfied; it is reslotted to the next valid window.
\end{itemize}

\subsubsection*{4. Operating Constraints}
\phantomsection

\begin{itemize}
\item Cadence: 6 sessions per week is mandatory. No two-a-days. No make-up doubling. Missed sessions remain missed and are logged accurately.

\item One-session-per-day rule:
  \begin{itemize}
  \item The daily training session is the single, unified daily training event closed by one unified log entry. Splitting required elements into separate events is non-compliant for evidence purposes unless executed and logged as one continuous session with a single unified entry and Cardio remains continuous.
  \end{itemize}

\item Session structure is mandatory for a session to be admissible:
  \begin{itemize}
  \item Safety Self-Check
  \item warm-up
  \item mobility
  \item main work
  \item Cardio
  \item cooldown
  \item recovery
  \item post-session notes
  \end{itemize}

\item Cardio occurs in every session and must meet the minimum standard: minimum 30 minutes continuous per session. If Cardio is under the minimum or not continuous, the session is \textbf{INVALID} and provides no gate credit.

\item Daily baseline activity:
  \begin{itemize}
  \item 10,000 steps/day baseline is recorded as a daily floor and feasibility monitor.
  \item Steps do not count as a session and do not satisfy session requirements.
  \end{itemize}

\item Tests are scheduled only inside Phase 00 Deload+Test windows, hard-locked: Weeks 4, 8, 12, 16. No informal tests in build weeks.

\item Phase 00 high-risk exclusions (non-negotiable):
  \begin{itemize}
  \item no ruck time trials,
  \item no run time trials beyond walk-based instruments defined in the Phase 00 test batteries,
  \item no swim time trials,
  \item no underwater,
  \item no maximal strength testing and no true \textbf{1RM} testing.
  \end{itemize}

\item Environmental safety overrides schedule and performance:
  \begin{itemize}
  \item unsafe conditions trigger substitution or reslot posture per Stage 3,
  \item compromised testing conditions are treated as \textbf{INVALID} for gate use.
  \end{itemize}

\item Compromised term handling:
  \begin{itemize}
  \item if source material uses “Compromised,” treat it as \textbf{INVALID} test evidence for gate decisions and as “not yet tested” until repeated in a valid window.
  \end{itemize}
\end{itemize}

\subsubsection*{5. Primary Adaptations and Physiological Focus}
\phantomsection

\begin{itemize}
\item Aerobic base via low-impact locomotion:
  \begin{itemize}
  \item establish tolerance for continuous, repeatable Cardio exposure without mechanics breakdown,
  \item stabilize pacing discipline so “hard days” do not appear by accident,
  \item improve recovery throughput by maintaining consistent low-cost aerobic work.
  \end{itemize}

\item Tissue tolerance and foot interface stability:
  \begin{itemize}
  \item reduce hotspots and prevent progression into blisters,
  \item maintain stable gait under longer walking durations,
  \item establish repeatable footwear and sock interface behavior.
  \end{itemize}

\item Neuromuscular re-education and movement competency:
  \begin{itemize}
  \item establish stable bracing and controlled joint ranges,
  \item improve movement quality for chair and incline variants under locked setup conditions,
  \item reduce compensatory mechanics that increase injury risk when volume increases later.
  \end{itemize}

\item Movement quality and mobility readiness:
  \begin{itemize}
  \item improve ankle dorsiflexion symmetry and shoulder flexion control under standardized checks,
  \item reduce pain-driven compensation and improve repeatability of screens.
  \end{itemize}

\item Recovery discipline and readiness transparency:
  \begin{itemize}
  \item stable sleep logging,
  \item soreness/fatigue tracking,
  \item pain mapping,
  \item session RPE stability to support safe progression decisions and to prevent fatigue contamination of test windows.
  \end{itemize}

\item \textbf{BODYCOMP} directionality and trend stabilization:
  \begin{itemize}
  \item establish clear trend directionality under stable method conditions,
  \item do not claim single-point success; trend logic is the evidence standard.
  \end{itemize}
\end{itemize}

\subsubsection*{6. Primary Modalities}
\phantomsection

\begin{itemize}
\item Low-impact Cardio modalities:
  \begin{itemize}
  \item walking (outdoor route) and/or treadmill walking under stable incline policy,
  \item walk-based monitor when selected as the locked aerobic monitor,
  \item indoor substitutes only when validity/logging preserves comparability,
  \item no running-specific constructs; impact escalation is prohibited in Phase 00.
  \end{itemize}

\item Calisthenics regressions and trunk endurance:
  \begin{itemize}
  \item incline push-up variant (\textbf{ME-UP-PU-INC} use Stage 2 \textbf{ID}),
  \item plank trunk endurance (\textbf{CAL-PLANK} use Stage 2 \textbf{ID}),
  \item any additional calisthenics tests (push-ups or sit-ups) are optional baseline-only and only attempted when strictly viable and risk-neutral.
  \end{itemize}

\item Lower-body function:
  \begin{itemize}
  \item chair sit-to-stand (\textbf{ME-L-STS-30} use Stage 2 \textbf{ID}) under locked chair height and arms policy.
  \end{itemize}

\item Mobility and movement screens:
  \begin{itemize}
  \item \textbf{MOB} Bundle (Requires Stage 8 review - no reconciled Stage 2 \textbf{ID} assigned), including ankle dorsiflexion and shoulder flexion wall posture plus movement screen.
  \end{itemize}

\item \textbf{DURABILITY} monitoring:
  \begin{itemize}
  \item pain flags, gait drift detection, foot hotspot monitoring, mechanics change detection.
  \end{itemize}

\item \textbf{RECOVERY} monitoring:
  \begin{itemize}
  \item sleep, soreness, fatigue, and session RPE trends used to downshift or reslot rather than push through.
  \end{itemize}
\end{itemize}

\begin{landscape}

\subsubsection*{7. Weekly Structure}
\phantomsection

Phase 00 is executed as a stable six-session rhythm with the same session-structure elements present every day. Weekly structure is intent-level only; it defines what the week is for and what it must not do (no impact escalation, no maximal testing, no unsafe catch-up). All sessions include Cardio, and Cardio must remain continuous and meet the minimum standard.

\paragraph*{Required Table 2 — Weekly Structure Table, Sessions 1-6}

\begingroup
\scriptsize
\setlength{\tabcolsep}{2pt}
\renewcommand{\arraystretch}{1.12}

\begin{longtblr}{
  width=\linewidth,
  rowhead=1,
  colspec = {
    Q[l,wd=0.70in]
    X[l,1.55]
    X[l,1.15]
    X[l,2.35]
    X[l,2.10]
  },
  hlines,
  vlines,
  cells = {hsep=6pt, vsep=6pt, halign=l, valign=t},
  row{odd} = {bg=white},
  row{even} = {bg=LightGray},
  row{1} = {bg=StageBlue!15, fg=black, font=\bfseries, halign=l, valign=t},
}
\Hdr{Session \#} &
\Hdr{Primary Focus} &
\Hdr{Main Work Domains} &
\Hdr{Cardio Modality/Intent} &
\Hdr{Notes/Risk Controls} \\

1 &
Low-risk aerobic consistency + movement prep &
\textbf{RECOVERY}; \textbf{DURABILITY} &
Cardio: walk or treadmill walk; minutes recorded (minimum standard applies); intensity anchor recorded as RPE 3–5 and talk-test cue “full sentences” to “short sentences”; continuity required &
Low-impact bias. If gait changes or hotspots occur, downshift immediately and document; no push through. Confirm footwear interface consistency. \\

2 &
Mobility emphasis + trunk endurance baseline &
\textbf{DURABILITY}; \textbf{CAL} &
Cardio: low-impact steady modality; minutes recorded (minimum standard applies); intensity anchor RPE 3–4 with talk-test “full sentences”; continuity required &
Technique-first posture. Any mechanics-altering pain triggers \textbf{STOP} \textbf{NOW} for that exposure and substitution to lowest risk. Maintain stable plank setup conditions for repeatability. \\

3 &
Aerobic tolerance + foot interface governance &
\textbf{DURABILITY}; \textbf{RECOVERY} &
Cardio: walking-focused exposure under stable footwear/route policy; minutes recorded (minimum standard applies); intensity anchor RPE 3–5 with talk-test; continuity required &
Foot inspection before and after session. Hotspot response is immediate and documented. Environment hazards trigger indoor substitution without reducing Cardio time. \\

4 &
Movement pattern control for chair and incline variants &
\textbf{DURABILITY}; \textbf{CAL} &
Cardio: low-impact steady modality; minutes recorded (minimum standard applies); intensity anchor RPE 3–5 with talk-test; continuity required &
Preserve lock-field readiness (chair height and incline height) without turning the session into a test outside protected windows. Avoid novelty in setup variables. \\

5 &
Recovery-biased aerobic throughput &
\textbf{RECOVERY}; \textbf{DURABILITY} &
Cardio: lowest-risk modality available; minutes recorded (minimum standard applies); intensity anchor RPE 2–4 with talk-test “full sentences”; continuity required &
Preserve cadence and recovery throughput. If readiness is poor, reduce main work dose and complexity; maintain Cardio minimum unless \textbf{STOP} \textbf{NOW} terminates. \\

6 &
Weekly consolidation + readiness check &
\textbf{RECOVERY}; \textbf{DURABILITY}; \textbf{CAL} (regressed) &
Cardio: walk-based or indoor substitute; minutes recorded (minimum standard applies); intensity anchor RPE 3–4 with talk-test; continuity required &
Consolidation, not a proving day. No novelty introduced. If entering a test window, treat this as stabilization support and ensure logs are complete. \\

\end{longtblr}

\endgroup

Weekly structure rules (binding):

\begin{itemize}
\item Cardio is present every session and must meet the minimum standard: minimum 30 minutes continuous. If Cardio is under the minimum or not continuous, the session is \textbf{INVALID} and provides no gate credit.
\item The weekly structure table defines intent only; no prescriptions (sets/reps/intervals) are authorized in this deliverable.
\item Session validity tagging is mandatory (\textbf{VALID} / \textbf{MODIFIED} / \textbf{INVALID}) and is driven by required element completion plus required logging completeness.
\item Phase 00 default intensity posture is conservative. Any drift into uncontrolled high intensity (RPE 8–10 patterns, repeated talk-test failure) is treated as a risk event requiring downshift and review.
\end{itemize}

\end{landscape}

\subsubsection*{8. Progression Rules}
\phantomsection

Progression rules exist to prevent “too much too soon,” prevent evidence contamination, and preserve adherence under obese novice posture.

Progression governors (must remain stable before advancing exposure):

\begin{itemize}
\item symptom stability:
  \begin{itemize}
  \item no \textbf{STOP} \textbf{NOW} triggers,
  \item no persistent dizziness, disproportionate dyspnea, or abnormal intolerance.
  \end{itemize}

\item mechanics stability:
  \begin{itemize}
  \item no gait-altering pain,
  \item no compensatory movement required to complete the exposure.
  \end{itemize}

\item foot integrity stability:
  \begin{itemize}
  \item hotspots addressed early,
  \item no progression into blisters or maceration patterns,
  \item no footwear-driven gait change.
  \end{itemize}

\item recovery stability:
  \begin{itemize}
  \item sleep opportunity and fatigue trends do not collapse in a way that predictably contaminates tests,
  \item RPE drift is interpreted conservatively.
  \end{itemize}

\item documentation completeness:
  \begin{itemize}
  \item missing fields are not “filled later,”
  \item same-day logs remain complete enough to audit validity.
  \end{itemize}
\end{itemize}

Progression levers (one at a time; conservative default):

\begin{itemize}
\item time-on-feet (increase only when mechanics and foot integrity are stable),
\item exposure continuity (maintain continuous Cardio rather than adding intensity spikes),
\item movement complexity (increase only when movement quality is stable),
\item testability (lock fields stabilized before a test is treated as admissible).
\end{itemize}

Lock-field stabilization rule (mandatory for testable items):

\begin{itemize}
\item Locked fields are setup variables that must remain consistent across Weeks 4–12 to preserve comparability.
  \begin{itemize}
  \item \textbf{ME-L-STS-30} lock fields: chair height; arms policy; timer method.
  \item \textbf{ME-UP-PU-INC} lock fields: incline height; hand/foot posture standard as defined by the protocol; timer/counting posture.
  \item \textbf{CAL-PLANK} lock fields: surface/mat posture; camera/timing posture if artifact evidence is required by Stage 2; start/stop criteria.
  \item \textbf{AER-WALK-015} lock fields: mode (outdoor route or treadmill), route/machine identity, footwear interface, surface class, timing method.
  \end{itemize}
\end{itemize}

No escalation by substitution:

\begin{itemize}
\item If a planned low-risk exposure cannot be performed, substitute to an equal-or-lower risk modality.
\item Substitution must not increase impact, load, or complexity.
\item Substitution is documented, and if it affects admissibility for a test window it is treated as \textbf{INVALID} for gate use.
\end{itemize}

No novelty near test windows:

\begin{itemize}
\item Weeks 4, 8, 12, and 16 are protected.
\item Avoid unlogged changes to:
  \begin{itemize}
  \item footwear interface,
  \item route/machine identity,
  \item timer method,
  \item chair height/incline height.
  \end{itemize}
\end{itemize}

If novelty is unavoidable, log it as a comparability break and treat the test attempt conservatively (\textbf{INVALID} for gate use unless Stage 2 validity conditions still support comparability).

\subsubsection*{9. Deload and Test Placement}
\phantomsection

\begin{itemize}
\item Hard-locked Deload+Test windows for Phase 00: Weeks 4, 8, 12, 16.
\end{itemize}

Purpose of Deload+Test windows in Phase 00:

\begin{itemize}
\item reduce fatigue interference and orthopedic risk so test outcomes can be \textbf{VALID} and comparable,
\item stabilize lock fields and environment/tool conditions,
\item concentrate evidence production into protected windows.
\end{itemize}

Build and intensify weeks:

\begin{itemize}
\item Weeks 1–3, 5–7, 9–11, 13–15 remain execution-focused.
\item Major test batteries are not scheduled in these weeks.
\item Any testing-like exposure in these weeks is treated as training information and is not used for gate outcomes.
\end{itemize}

Deload+Test week conversion posture:

\begin{itemize}
\item If a Deload+Test week cannot support \textbf{VALID} conditions due to illness, pain drift, missing lock fields, environmental uncertainty, or tool drift:
  \begin{itemize}
  \item convert the week to deload-only or check-in-only,
  \item reslot maximal tests to the next protected window.
  \end{itemize}
\end{itemize}

\newpage
\subsubsection*{10. Test Battery}
\phantomsection

Gate-grade tests occur only in Weeks 4, 8, 12, 16. Phase 00 uses conservative batteries that emphasize movement readiness, trunk endurance, and walk-based aerobic monitors. Underwater is not scheduled.

Gate-admissible versus baseline/trend-only rule (binding):

\begin{itemize}
\item Gate-admissible evidence for Phase 00 exit is produced in Week 16 only, and only if test validity is \textbf{VALID}.
\item Week 4, 8, and 12 test events are baseline/trend or confirmation events supporting comparability and readiness monitoring; they do not, by default, authorize Phase progression.
\end{itemize}

\textbf{CAL} artifact posture (Phase 00 relevance):

\begin{itemize}
\item Any \textbf{CAL} test included in a gate-grade decision window must satisfy the Stage 2.C \textbf{CAL} artifact evidence standard when the protocol requires it.
\item Missing/non-compliant artifact evidence renders the \textbf{CAL} result \textbf{INVALID} for gate use even if the numeric result appears strong.
\end{itemize}

\begin{landscape}

\paragraph*{Required Table 3 — Test Battery Table}

\begingroup
\scriptsize
\setlength{\tabcolsep}{2pt}
\renewcommand{\arraystretch}{1.12}

\begin{longtblr}{
  width=\linewidth,
  rowhead=1,
  colspec = {
    X[l,1.55]
    X[l,1.55]
    Q[l,wd=1.05in]
    X[l,2.25]
    X[l,1.60]
  },
  hlines,
  vlines,
  cells = {hsep=6pt, vsep=6pt, halign=l, valign=t},
  row{odd} = {bg=white},
  row{even} = {bg=LightGray},
  row{1} = {bg=StageBlue!15, fg=black, font=\bfseries, halign=l, valign=t},
}
\Hdr{Test Timing} &
\Hdr{Test \textbf{ID}/Name} &
\Hdr{Domain} &
\Hdr{Validity Conditions} &
\Hdr{Pass Threshold} \\

Week 4 — Baseline Battery, Deload+Test &
\textbf{SMB} (Requires Stage 8 review - no reconciled Stage 2 \textbf{ID} assigned) &
\textbf{DURABILITY} &
Safety screen passed; no red-flag symptoms; no mechanics-altering pain; lock fields and environment/tool conditions fully logged; Compromised treated as \textbf{INVALID} for gate use &
Baseline capture only; gate use occurs at Week 16 only if \textbf{VALID} \\

Week 4 — Baseline Battery, Deload+Test &
\textbf{MOB} Bundle: \textbf{MOB-MOV-SCREEN} (Requires Stage 8 review - no reconciled Stage 2 \textbf{ID} assigned) &
\textbf{DURABILITY} &
Standard movement screen format held stable; pain and compensation noted; required fields complete; Compromised treated as \textbf{INVALID} for gate use &
Pass or only minor flags with documented plan; no pain-driven compensation (gate use at Week 16 only if \textbf{VALID}) \\

Week 4 — Baseline Battery, Deload+Test &
\textbf{MOB} Bundle: \textbf{MOB-ANKDF-K2W} (Requires Stage 8 review - no reconciled Stage 2 \textbf{ID} assigned) &
\textbf{DURABILITY} &
Knee-to-wall method stable; side-specific measurement; pain noted; required fields complete &
At least 3.0 in each side; asymmetry at most 0.5 in; no pain (gate use at Week 16 only if \textbf{VALID}) \\

Week 4 — Baseline Battery, Deload+Test &
\textbf{MOB} Bundle: \textbf{MOB-SHFLEX-WALL} (Requires Stage 8 review - no reconciled Stage 2 \textbf{ID} assigned) &
\textbf{DURABILITY} &
Wall test method stable; rib flare compensation checked; pain noted; required fields complete &
Pass or gap at most 2.0 in without rib flare compensation or pain (gate use at Week 16 only if \textbf{VALID}) \\

Week 4 — Baseline Battery, Deload+Test &
\textbf{MOB} Bundle: \textbf{MOB-STS-QUAL} (Requires Stage 8 review - no reconciled Stage 2 \textbf{ID} assigned) &
\textbf{DURABILITY} &
Standard observation criteria stable; compensation and pain noted; required fields complete &
Baseline capture only unless Stage 2 defines a pass criterion; gate use at Week 16 only if \textbf{VALID} \\

Week 4 — Baseline Battery, Deload+Test &
\textbf{ME-L-STS-30} (Requires Stage 8 review - no reconciled Stage 2 \textbf{ID} assigned) &
\textbf{DURABILITY} &
Chair height locked and recorded; arms policy locked and recorded; timer method stable; same-day logging; no pain-driven compensation &
At least 12 reps in 30 s (gate use at Week 16 only if \textbf{VALID}) \\

Week 4 — Baseline Battery, Deload+Test &
\textbf{ME-UP-PU-INC} (Requires Stage 8 review - no reconciled Stage 2 \textbf{ID} assigned) &
\textbf{CAL} &
Incline height locked and recorded; rep standard held stable; same-day logging; no mechanics-altering pain &
At least 15 reps at fixed incline height (gate use at Week 16 only if \textbf{VALID}) \\

Week 4 — Baseline Battery, Deload+Test &
\textbf{CAL-PLANK} (Requires Stage 8 review - no reconciled Stage 2 \textbf{ID} assigned) &
\textbf{CAL} &
Standard plank position; termination on form failure; timer method stable; same-day logging &
At least 1:00 strict (gate use at Week 16 only if \textbf{VALID}) \\

Week 4 — Baseline Battery, Deload+Test &
\textbf{AER-STEP-3MIN} (choose one locked aerobic monitor) (Requires Stage 8 review - no reconciled Stage 2 \textbf{ID} assigned) &
\textbf{AER} &
Locked selection must be declared and held stable through Week 12; mode and setup logged; environment/tool stable; same-day logging &
Trend/baseline-only; no Phase 00 advancement claim based on this metric \\

Week 4 — Baseline Battery, Deload+Test &
\textbf{AER-WALK-30} (choose one locked aerobic monitor) (Requires Stage 8 review - no reconciled Stage 2 \textbf{ID} assigned) &
\textbf{AER} &
Locked selection must be declared and held stable through Week 12; footwear and route/surface class logged; mode logged; same-day logging &
Trend/baseline-only; no Phase 00 advancement claim based on this metric \\

Week 8 — Mid Check, Deload+Test &
\textbf{SMB} (Requires Stage 8 review - no reconciled Stage 2 \textbf{ID} assigned) &
\textbf{DURABILITY} &
Same \textbf{SMB} validity posture; full logging; Compromised treated as \textbf{INVALID} for gate use &
Mid-phase check-in only; gate use at Week 16 only if \textbf{VALID} \\

Week 8 — Mid Check, Deload+Test &
\textbf{MOB} Bundle (all components above) (use Stage 2 \textbf{ID}) &
\textbf{DURABILITY} &
All lock fields held stable; pain and compensation documented; required fields complete &
Mid-phase check-in only; gate standards apply at Week 16 only if \textbf{VALID} \\

Week 8 — Mid Check, Deload+Test &
\textbf{ME-L-STS-30} (Requires Stage 8 review - no reconciled Stage 2 \textbf{ID} assigned) &
\textbf{DURABILITY} &
Chair height locked; arms policy locked; timer method stable &
Mid-phase check-in only; gate threshold applies at Week 16 only if \textbf{VALID} \\

Week 8 — Mid Check, Deload+Test &
\textbf{ME-UP-PU-INC} (Requires Stage 8 review - no reconciled Stage 2 \textbf{ID} assigned) &
\textbf{CAL} &
Incline height locked; rep standard stable &
Mid-phase check-in only; gate threshold applies at Week 16 only if \textbf{VALID} \\

Week 8 — Mid Check, Deload+Test &
\textbf{CAL-PLANK} (Requires Stage 8 review - no reconciled Stage 2 \textbf{ID} assigned) &
\textbf{CAL} &
Standard position; timer method stable &
Mid-phase check-in only; gate threshold applies at Week 16 only if \textbf{VALID} \\

Week 8 — Mid Check, Deload+Test &
\textbf{AER-WALK-015}, mode logged (use Stage 2 \textbf{ID}) &
\textbf{RUN} &
Mode logged; footwear and route/surface class logged; stable gait required; same-day logging; Compromised treated as \textbf{INVALID} for gate use &
Gate threshold applies at Week 16 only if \textbf{VALID} \\

Week 12 — Confirm, Deload+Test &
Repeat Week 4 battery (all items above) (use Stage 2 \textbf{ID}) &
\textbf{DURABILITY} / \textbf{CAL} / \textbf{RUN} &
Repeat under the same lock fields; locked aerobic monitor selection must match Week 4 choice; Compromised treated as \textbf{INVALID} for gate use &
Confirmation only; gate thresholds apply at Week 16 only if \textbf{VALID} \\

Week 16 — Exit Gate, Deload+Test &
\textbf{SMB} (Requires Stage 8 review - no reconciled Stage 2 \textbf{ID} assigned) &
\textbf{DURABILITY} &
Full lock fields and logging; no red flags; no mechanics-altering pain; Compromised treated as \textbf{INVALID} for gate use &
\textbf{PASS} required for gate consideration (\textbf{SMB} clearance) \\

Week 16 — Exit Gate, Deload+Test &
\textbf{MOB} Bundle (required gate components) (use Stage 2 \textbf{ID}) &
\textbf{DURABILITY} &
All gate-relevant components executed under stable method; pain must be absent for required elements &
\textbf{MOB-ANKDF-K2W}, \textbf{MOB-SHFLEX-WALL}, \textbf{MOB-MOV-SCREEN} must meet the Phase 00 exit standards \\

Week 16 — Exit Gate, Deload+Test &
\textbf{ME-L-STS-30} (Requires Stage 8 review - no reconciled Stage 2 \textbf{ID} assigned) &
\textbf{DURABILITY} &
Chair height locked; arms policy locked; timer stable; no pain-driven compensation &
At least 12 reps in 30 s \\

Week 16 — Exit Gate, Deload+Test &
\textbf{ME-UP-PU-INC} (Requires Stage 8 review - no reconciled Stage 2 \textbf{ID} assigned) &
\textbf{CAL} &
Incline height locked; rep standard stable; no mechanics-altering pain &
At least 15 reps at fixed incline height \\

Week 16 — Exit Gate, Deload+Test &
\textbf{CAL-PLANK} (Requires Stage 8 review - no reconciled Stage 2 \textbf{ID} assigned) &
\textbf{CAL} &
Standard plank position; termination on form failure; timer stable &
At least 1:00 strict \\

Week 16 — Exit Gate, Deload+Test &
\textbf{AER-WALK-015}, mode logged (use Stage 2 \textbf{ID}) &
\textbf{RUN} &
Mode logged; stable gait; no symptom flags; route/footwear/mode lock fields held stable for comparability &
At most 25:00 \\

Week 16 — Exit Gate, Deload+Test &
\textbf{REC-ADH-WK} (adherence audit) (use Stage 2 \textbf{ID}) &
\textbf{RECOVERY} &
Counting rules applied; admissible session definition applied; logs complete &
At least 85\% across Phase 00; at least 82 of 96 sessions \\

Week 16 — Exit Gate, Deload+Test &
\textbf{BODYCOMP} trend directionality evidence (\textbf{BC-WT} and/or \textbf{BC-CIRC-WAIST}) (use Stage 2 \textbf{ID}) &
\textbf{BODYCOMP} &
Same method and cadence rules; no backfilled entries; trend is interpretable &
Clear downward trend vs Phase start; no sustained upward drift without documented cause \\

Week 16 — Optional baseline-only, attempt only if strictly viable &
\textbf{CAL-PU2} optional (use Stage 2 \textbf{ID}) &
\textbf{CAL} &
Only if it does not increase risk; full logging; stop posture enforced &
Baseline-only; not used to \textbf{ADVANCE} out of Phase 00 \\

Week 16 — Optional baseline-only, attempt only if strictly viable &
\textbf{CAL-CU2} (optional; Requires Stage 8 review - no reconciled Stage 2 \textbf{ID} assigned) &
\textbf{CAL} &
Only if it does not increase risk; full logging; stop posture enforced &
Baseline-only; not used to \textbf{ADVANCE} out of Phase 00 \\

\end{longtblr}

\endgroup

\end{landscape}

Phase 00 test lock-field register (required for comparability; stored in the test record Notes/\textbf{MEASURE} fields):

\begin{itemize}
\item \textbf{ME-L-STS-30}:
  \begin{itemize}
  \item chair height (measured or described consistently),
  \item arms policy (arms crossed / hands on thighs / other fixed policy),
  \item timer method identity.
  \end{itemize}

\item \textbf{ME-UP-PU-INC}:
  \begin{itemize}
  \item incline height (measured or described consistently),
  \item hand position standard (if the protocol defines),
  \item rep counting authority posture.
  \end{itemize}

\item \textbf{CAL-PLANK}:
  \begin{itemize}
  \item surface/mat posture,
  \item start/stop criteria,
  \item timing method identity,
  \item artifact compliance posture if required by Stage 2.
  \end{itemize}

\item \textbf{AER-WALK-015}:
  \begin{itemize}
  \item mode: route or treadmill,
  \item Environment Unit \textbf{ID} (\textbf{ROUTE}-2E-### or \textbf{TM}-2E-###) if applicable under Stage 2 logging,
  \item footwear interface label,
  \item timing method identity.
  \end{itemize}
\end{itemize}

\subsubsection*{11. Exit Standards}
\phantomsection

Phase 00 exit decisions are made at Week 16 only, using gate-admissible evidence that is \textbf{VALID}. If evidence is \textbf{INVALID}, it is treated as not yet tested and cannot support \textbf{ADVANCE}.

Minimum required exit standards (all are binding; thresholds are the ones listed here):

\begin{enumerate}
\item Adherence and execution discipline

\begin{itemize}
\item \textbf{REC-ADH-WK} (Requires Stage 8 review - no reconciled Stage 2 \textbf{ID} assigned):
  \begin{itemize}
  \item at least 85\% across Phase 00,
  \item at least 82 of 96 sessions,
  \item counting rules (operational definition for this phase row):
    \begin{itemize}
    \item denominator is 96 scheduled sessions (16 weeks × 6 sessions/week),
    \item numerator counts sessions that are admissible for adherence credit:
      \begin{itemize}
      \item \textbf{VALID} sessions count,
      \item \textbf{MODIFIED} sessions count only when admissible under Stage 1 governance posture and fully documented,
      \item \textbf{INVALID} sessions and missed sessions do not count,
      \end{itemize}
    \item any session with Cardio under the minimum standard or not continuous is \textbf{INVALID} and cannot be counted.
    \end{itemize}
  \end{itemize}
\end{itemize}

\item Aerobic monitor gate item

\begin{itemize}
\item \textbf{AER-WALK-015} (Requires Stage 8 review - no reconciled Stage 2 \textbf{ID} assigned) at most 25 minutes:
  \begin{itemize}
  \item mode logged,
  \item stable gait,
  \item no symptom flags,
  \item route/footwear/mode lock fields held stable for comparability.
  \end{itemize}
\end{itemize}

\item Trunk endurance

\begin{itemize}
\item \textbf{CAL-PLANK} (use Stage 2 \textbf{ID}) at least 1 minute strict:
  \begin{itemize}
  \item strict position and termination on form failure,
  \item timing posture verifiable.
  \end{itemize}
\end{itemize}

\item Lower-body function and tolerance

\begin{itemize}
\item \textbf{ME-L-STS-30} (use Stage 2 \textbf{ID}) at least 12 reps in 30 seconds:
  \begin{itemize}
  \item chair height locked,
  \item arms policy locked,
  \item no pain-driven compensation.
  \end{itemize}
\end{itemize}

\item Upper-body function (regressed standard)

\begin{itemize}
\item \textbf{ME-UP-PU-INC} (use Stage 2 \textbf{ID}) at least 15 reps at fixed incline height:
  \begin{itemize}
  \item incline height locked,
  \item rep standard stable,
  \item no mechanics-altering pain.
  \end{itemize}
\end{itemize}

\item Mobility gate elements

\begin{itemize}
\item \textbf{MOB-ANKDF-K2W} (use Stage 2 \textbf{ID}):
  \begin{itemize}
  \item at least 3.0 inches each side,
  \item asymmetry at most 0.5 inch,
  \item no pain.
  \end{itemize}

\item \textbf{MOB-SHFLEX-WALL} (use Stage 2 \textbf{ID}):
  \begin{itemize}
  \item pass or gap at most 2.0 inches,
  \item without rib flare compensation,
  \item no pain.
  \end{itemize}

\item \textbf{MOB-MOV-SCREEN} (use Stage 2 \textbf{ID}):
  \begin{itemize}
  \item pass or only minor flags with documented plan,
  \item no pain-driven compensation.
  \end{itemize}
\end{itemize}

\item \textbf{DURABILITY} stability proof (trend and documentation-based)

\begin{itemize}
\item No hotspots progressing into blisters.
\item No gait drift under longer walking.
\item No repeated mechanics-altering pain episodes.
\item No repeated \textbf{STOP} \textbf{NOW} symptom categories.
\item No repeated sleep collapse patterns that produce RPE drift spikes and session invalidity.
\item Must be demonstrated via logs across Phase 00, not by a single day.
\end{itemize}

\item \textbf{BODYCOMP} directionality (trend evidence only)

\begin{itemize}
\item \textbf{BC-WT} or \textbf{BC-CIRC-WAIST} (Requires Stage 8 review - no reconciled Stage 2 \textbf{ID} assigned) shows a clear downward trend vs Phase start under the same method and cadence rules.
\item No sustained upward drift without documented cause (access disruption, illness, travel, measurement method change recorded as comparability break).
\end{itemize}
\end{enumerate}

If not met disposition (mandatory):

\begin{itemize}
\item Do not advance to Phase 01.
\item Mandatory disposition: \textbf{HOLD} Phase 01 entry, extend Phase 00 only under documented control, then retest in the next Deload+Test window.
\item Retest only what is required to clear the failed standard; do not stack unnecessary tests.
\end{itemize}

\subsubsection*{12. Risk Controls and Stop Rules}
\phantomsection

Phase 00 risk posture is conservative and explicitly biased toward preventing \textbf{MSK} overload, cardiometabolic events, and evidence contamination.

Stop posture (governance-level; applies to sessions and tests):

\begin{itemize}
\item Red-flag symptoms abort testing immediately and invoke \textbf{STOP} \textbf{NOW} posture.
\item Mechanics-altering pain (gait change, limping, compensatory movement) terminates the provoking exposure and triggers substitution.
\item Acute injury with instability or swelling triggers \textbf{STOP} \textbf{NOW} and escalation posture per governance.
\end{itemize}

Phase 00 domain restrictions (non-negotiable):

\begin{itemize}
\item No impact running time trials (only walk-based instruments).
\item No ruck time trials and no load carriage testing.
\item No swim time trials.
\item No underwater or breath-hold exposure under any circumstances.
\item No maximal strength testing.
\end{itemize}

Foot integrity governance:

\begin{itemize}
\item Foot hotspots are early-warning triggers; correct action is treat-and-adjust, not push-through.
\item Foot hotspot presence in the lead-in to a test window biases toward reslotting walk-based tests if gait stability cannot be preserved.
\item Footwear interface is treated as a controlled input; footwear changes near Weeks 4/8/12/16 are comparability breaks.
\end{itemize}

Environment safety governance:

\begin{itemize}
\item Unsafe footing, poor visibility, storms, unsafe heat, or poor air quality triggers indoor substitution or reslot.
\item Gate-grade testing is not attempted under environment hazard conditions; compromised attempts are \textbf{INVALID} and reslotted.
\end{itemize}

\textbf{INVALID} handling:

\begin{itemize}
\item Any test that cannot meet \textbf{VALID} conditions due to illness, pain drift, missing lock fields, environmental uncertainty, or equipment change is tagged Compromised/\textbf{INVALID} and treated as not yet tested for gate decisions until repeated in the next valid window.
\end{itemize}

\subsubsection*{13. Required Capabilities and Dependencies}
\phantomsection

Phase 00 requires minimal, function-linked capabilities. No shopping lists are permitted; dependencies are expressed as required capabilities with owners, due-by timing, and fail actions.

\begin{landscape}

\paragraph*{Required Table 4 — Dependencies Table, Phase 00 Minimal Prerequisites}

\begingroup
\scriptsize
\setlength{\tabcolsep}{2pt}
\renewcommand{\arraystretch}{1.12}

\begin{longtblr}{
  width=\linewidth,
  rowhead=1,
  colspec = {
    Q[l,wd=1.05in]
    X[l,2.45]
    Q[l,wd=1.25in]
    Q[l,wd=1.10in]
    X[l,2.55]
  },
  hlines,
  vlines,
  cells = {hsep=6pt, vsep=6pt, halign=l, valign=t},
  row{odd} = {bg=white},
  row{even} = {bg=LightGray},
  row{1} = {bg=StageBlue!15, fg=black, font=\bfseries, halign=l, valign=t},
}
\Hdr{Dependency \textbf{ID}} &
\Hdr{Required Capability} &
\Hdr{Due-By} &
\Hdr{Owner} &
\Hdr{Fail Action} \\

\textbf{DEP-1C-001} &
Medical clearance posture for starting Phase 00 (documented restrictions become binding) &
Before Phase 00 Week 1 &
Medical &
\textbf{MEDICAL} \textbf{HOLD} posture for Phase start; do not execute Phase 00 as written until clearance exists; tests in Weeks 4/8/12/16 are not gate-admissible until clearance posture is satisfied \\

\textbf{DEP-1C-002} &
Stop-rule understanding and escalation access is operational &
Before Phase 00 Week 1 &
Coach Lead &
Cap intensity and complexity; restrict to lowest-risk modalities; treat any red-flag symptom as immediate termination; \textbf{HOLD} posture for any advancement claims until validated \\

\textbf{DEP-1C-003} &
Safe primary Cardio route or indoor substitute modality available &
Before Phase 00 Week 1 &
Trainee &
If no safe route or substitute exists, sessions may become \textbf{INVALID}; default \textbf{HOLD} until a safe option exists; reslot any walk-based test items until environment feasibility is validated \\

\textbf{DEP-1C-004} &
Weather and environment substitution plan exists and is used &
By Phase 00 Week 2 &
Coach Lead &
Default to indoor substitutions; outdoor exposures restricted to safe conditions only; any test attempted under unsafe/unlogged conditions is \textbf{INVALID}; \textbf{HOLD} posture until plan exists \\

\textbf{DEP-1C-005} &
Logging method and daily/session compliance capability &
By Phase 00 Week 3 &
Trainee &
No advancement posture; \textbf{HOLD} by default; repeat 14-day log proof until pass; tests may be executed but are inadmissible for gate use without complete logs \\

\textbf{DEP-1C-006} &
Timing and measurement tools sufficient for validity (timer + steps method + weight method) &
By Phase 00 Week 2 &
Equipment Lead &
Any time-based test is \textbf{INVALID} without timing identity; postpone gate-grade tests; \textbf{HOLD} posture until tools exist and are stable \\

\textbf{DEP-1C-007} &
Footwear interface established and fit verified &
Before Phase 00 Week 1 &
Equipment Lead &
Cap impact and duration; prohibit any escalation; treat walk-based test items as \textbf{INVALID} until footwear stability is demonstrated without hotspots or gait change \\

\textbf{DEP-1C-008} &
Footcare supplies and hotspot control protocol available &
By Phase 00 Week 2 &
Equipment Lead &
Downshift exposure; substitute to lower-friction modalities; treat hotspot with mechanics change as immediate stop trigger; \textbf{HOLD} posture until stable foot integrity is demonstrated \\

\textbf{DEP-1C-011} &
Storage and maintenance minimum (drying, safe training lane) &
By Phase 00 Week 4 &
Trainee &
Reduce equipment footprint; prohibit unstable setups; restrict wet-environment exposures; \textbf{HOLD} posture if repeated breakdown prevents session validity \\

\textbf{DEP-1C-014} &
Schedule stability for 6 sessions/week (time blocks protected) &
By Phase 00 Week 2 &
Trainee &
Implement minimum viable scheduling; prohibit make-up volume; \textbf{HOLD} posture for any escalation until cadence proof restored \\

\textbf{DEP-1C-015} &
Binder access and version visibility (single source of truth) &
Before Phase 00 Week 1 &
Coach Lead &
Freeze change behavior; prohibit informal edits; restrict decisions to re-established doctrine; \textbf{HOLD} posture for advancement until access restored \\

\textbf{DEP-1C-016} &
Hydration access and container availability &
Before Phase 00 Week 1 &
Trainee &
Downshift intensity; avoid heat exposure; prefer indoor climate-controlled modality; terminate under \textbf{STOP} \textbf{NOW} posture if dizziness/heat signs appear; \textbf{HOLD} posture if persistent \\

\end{longtblr}

\endgroup

\end{landscape}

\subsubsection*{14. Validity Requirements}
\phantomsection

Validity requirements are enforceable controls; they determine admissibility and gate credit posture.

Session validity (Stage 1.F controls; Phase 00 enforcement posture):

\begin{itemize}
\item Session is \textbf{VALID} only if all required elements are completed and logged, including Cardio and recovery/post-session notes.
\item Cardio is minimum 30 minutes continuous per session. If Cardio is under the minimum or not continuous, the session is \textbf{INVALID} and provides no gate credit.
\item \textbf{MODIFIED} and \textbf{INVALID} sessions require reason-code posture (\textbf{SESSION-RC} per Stage 1.F) and a plain-language statement of what changed and why.
\end{itemize}

Test validity (Phase 00 batteries; Stage 1.F controls; Phase 00 Compromised mapping):

\begin{itemize}
\item Tests are \textbf{VALID} or \textbf{INVALID} only. “Compromised” is treated as \textbf{INVALID} for gate use.
\item If a scheduled test cannot meet \textbf{VALID} conditions due to illness, pain drift, missing lock fields, environmental uncertainty, or equipment change:
  \begin{itemize}
  \item tag Compromised/\textbf{INVALID},
  \item treat as not yet tested for gate decisions,
  \item repeat in the next valid window.
  \end{itemize}
\end{itemize}

Required minimum logging posture for Phase 00 test events (protocol-agnostic; aligns to Stage 2 tokens):

\begin{itemize}
\item \textbf{ENV}: environment/site or unit identity sufficient to interpret comparability (route/machine/site).
\item \textbf{EQUIP}: tool identity posture (timer method; chair or incline identifier; footwear interface descriptor).
\item \textbf{WARMUP}: completion recorded where applicable (or Not applicable when permitted).
\item \textbf{SAFETY}: pass/fail screen plus any stop triggers.
\item \textbf{MEASURE}: lock fields recorded (chair height; incline height; mode; arms policy; timing method) and any deviations flagged.
\end{itemize}

Evidence Ref posture:

\begin{itemize}
\item Evidence Ref is mandatory only when required by the governing Stage 2 protocol or logging rule. Missing Evidence Ref yields inadmissible posture for gate use and defaults to \textbf{INVALID} for gate use.
\end{itemize}

Gate admissibility posture for Phase 00 exit:

\begin{itemize}
\item Week 16 gate outcomes are based on gate-admissible evidence only.
\item Trend evidence (\textbf{BODYCOMP} directionality; \textbf{DURABILITY} stability) is required as documentation-based stability, not as a single-point claim.
\end{itemize}

\subsubsection*{15. Change Control Notes}
\phantomsection

\begin{itemize}
\item Phase 00 Deload+Test windows are hard-locked at Weeks 4, 8, 12, 16. Changing these windows requires patch control consistent with Stage 1 change control posture.

\item Allowed without patch (documentation required):
  \begin{itemize}
  \item swap sequencing of tests within a Deload+Test week while preserving validity and not introducing new tests,
  \item reslot a Compromised/\textbf{INVALID} test to the next Deload+Test window within Phase 00,
  \item downgrade a Deload+Test week to deload-only/check-in-only when safety, environment, or readiness makes gate-grade attempts inadmissible,
  \item remove any aquatic or underwater items (if ever mistakenly scheduled) consistent with safety posture (Phase 00 underwater is prohibited).
  \end{itemize}

\item Requires patch:
  \begin{itemize}
  \item add a new test week,
  \item extend phase duration beyond the defined extension disposition rule,
  \item change which battery items are treated as gate-admissible,
  \item change lock fields or pass thresholds.
  \end{itemize}

\item Any patch must align to Stage 1 change control posture and must not introduce new tests, new domains, or new thresholds.
\end{itemize}

\subsubsection*{16. Compliance Checklist}
\phantomsection

Phase 00 is compliant only when the following are true:

\begin{itemize}
\item Cadence:
  \begin{itemize}
  \item 6 sessions per week executed, with no hidden rest days,
  \item no two-a-days and no make-up doubling.
  \end{itemize}

\item Session structure and Cardio:
  \begin{itemize}
  \item every session includes required elements in order,
  \item Cardio is completed every session and meets the minimum standard (minimum 30 minutes continuous),
  \item any deficiency is tagged \textbf{MODIFIED} or \textbf{INVALID} with documentation.
  \end{itemize}

\item Logging and evidence:
  \begin{itemize}
  \item same-day logs exist for all sessions,
  \item recovery and post-session notes are complete enough to determine admissibility,
  \item missing fields are recorded as \textbf{MISSING} or \textbf{UNKNOWN} and are treated as inadmissible for gate decisions.
  \end{itemize}

\item Phase 00 test execution:
  \begin{itemize}
  \item batteries occur only in Weeks 4, 8, 12, 16,
  \item lock fields are recorded and held stable where required,
  \item any \textbf{INVALID} test is treated as not yet tested for gate decisions and reslotted.
  \end{itemize}

\item Safety posture:
  \begin{itemize}
  \item \textbf{STOP} \textbf{NOW} triggers are honored immediately,
  \item pain that alters mechanics terminates the exposure and triggers substitution/reslot posture,
  \item environment hazards trigger indoor substitution only when comparability can be preserved; otherwise reslot.
  \end{itemize}

\item Water governance:
  \begin{itemize}
  \item underwater is not scheduled and not attempted in Phase 00 under any circumstance,
  \item no breath-hold exposure is scheduled or attempted in Phase 00.
  \end{itemize}

\item Exit gate posture:
  \begin{itemize}
  \item \textbf{ADVANCE} to Phase 01 is permitted only when Week 16 exit standards are met with \textbf{VALID} evidence and required trend stability proof,
  \item if not met: \textbf{HOLD} Phase 01 entry, apply documented extension or repeat posture, and retest only the failed standards in the next Deload+Test window.
  \end{itemize}
\end{itemize}

\end{document}